\section{Discussion}
\label{sec:discussion}

This study examines what connectivity-based autism classification measures in ABIDE-I, how its probabilities and explanations should be interpreted, and which analysis choices influence reported performance. Five main contributions emerge. First, the ensemble probabilities show a small average overprediction of autism, and recalibration fitted only within the training partitions removes this offset without changing discrimination. Second, the scores of the linear ensemble members can be decomposed exactly into functional-system contributions, and these patterns reproduce across non-overlapping halves of the cohort at the system level but are less stable at the level of individual connections. Third, lower somatomotor connectivity in autistic participants persists and becomes stronger after within-site matching on recorded mean FD, while default-mode pairs account for the largest overall shares of classifier separation. Fourth, discrimination varies across cognitive-ability subgroups in this cohort, whereas classifier scores are not associated with symptom severity within autistic participants after correction. Fifth, the evaluated linear classifiers outperform the fixed-configuration tree ensembles on the same tangent-space features, while the small apparent advantage of windowed connectivity is not retained when covariance shrinkage is matched. Together, these findings show that reported ABIDE-I performance depends partly on cohort composition and analytical design. In particular, differences in cognitive composition can change the measured AUROC, which limits direct comparison between differently composed samples. The persistence of the somatomotor reduction after motion matching suggests that recorded mean FD alone does not explain the effect, although this exploratory result does not establish a biomarker or exclude residual motion. The classifier comparisons further support linear models as effective choices for these high-dimensional connectivity features, while the symptom analysis shows that diagnostic-group separation should not be interpreted as measurement of autism severity. The shrinkage analysis also shows that no benefit of temporal windowing is detected once covariance shrinkage is matched, so an apparent benefit can reflect regularisation instead of additional temporal information. Finally, the calibration and attribution analyses show that implementation choices can influence probability estimates and that interpretation is more stable at the level of functional systems than at the level of individual connections.

\subsection{Interpretation of the results}
\label{subsec:discussion-results}

Discrimination is maintained when an entire acquisition site is withheld from training. The primary ensemble reaches a participant-pooled AUROC of 0.776 under pooled ten-fold validation and 0.774 under LOSO validation (Table~\ref{tab:performance}). The similarity of these values indicates no clear loss in discrimination when the test site is excluded from model fitting. Across the 20 sites, the random-effects mean AUROC is 0.803, with a descriptive interval of 0.725--0.882. This mean weights sites by inverse variance, which favours sites with higher AUROC, so it exceeds the unweighted site mean (0.773) and the participant-pooled AUROC (0.774) and should not be compared with them as an estimate of the same quantity. Because the site estimates are based on overlapping training data, this interval is not interpreted as a validated prediction interval for a new centre. Site-level AUROC nevertheless ranges from 0.589 to 0.912, and five site intervals include 0.5, which reflects substantial uncertainty in sites containing only about 30 participants. Performance is also comparable with previously reported ABIDE-I results. In the wider cohort, the ensemble reaches an accuracy of 0.706 (95\% Wilson interval 0.678--0.733), compared with published values of 0.700, 0.703, and 0.708 \citep{Heinsfeldetal2018,Eslamietal2019,AlmuqhimSaeed2021}. These values support comparable performance rather than superiority. 

Cognitive composition, however, affects measured discrimination. In the unmatched comparison against the same 468 non-autistic participants, AUROC decreases from 0.873 for autistic participants with full-scale IQ below 85 to 0.719 for those with IQ 115 or above. Within autistic participants, classifier score is negatively associated with full-scale IQ ($\rho=-0.18$, $p=0.0004$), whereas the association is small in non-autistic participants ($\rho=0.05$, $p=0.32$). After IQ matching, discrimination remains highest in the lowest quartile but no longer declines consistently across the upper three, indicating that part of the original gradient reflects between-group IQ differences. The diagnosis-by-IQ interaction is nevertheless significant for full-scale and verbal IQ, so the evidence supports a stronger negative score--IQ association in autistic participants and variation of measured performance with cognitive composition, not a simple monotonic effect of IQ. The finding is also consistent with the verbal-IQ association previously reported in ABIDE \citep{Nielsenetal2013}, but independent validation is needed before an IQ-specific reliability claim can be made. Motion matching similarly changes the interpretation of group differences. Mean FD is positively associated with mean connection strength in this pipeline ($\rho=0.51$ across participants and $0.45$ within site) and with every within-somatomotor connection (mean $\rho=0.18$ after removing site), and the shift that the group difference in mean FD would produce in these connections through this linear relation is $+0.006$, whereas the observed site-adjusted group difference is $-0.046$. The somatomotor reduction is therefore not readily explained by the recorded relation between mean FD and connectivity, which does not exclude other motion effects. It also becomes stronger after within-site matching on recorded mean FD: under Freedman--Lane permutation, the number of reduced system pairs rises from 14 to 27 of 36 in 256 matched pairs (from 10 to 26 under the within-site permutation). The regional follow-up is descriptive because the system was chosen after inspection of the system-level results: the mean effect across the 21 somatomotor regions changes from $-0.215$ in all 863 participants to $-0.249$ in 280 matched pairs, and every regional effect is more negative in the matched sample. The direction agrees with earlier reports \citep{Chenetal2015,Holigaetal2019,Ilioskaetal2023}, although overlap with ABIDE prevents treating this agreement as independent replication. The feature-selection analysis further shows that reported performance depends strongly on how features are chosen. Selecting 2,000 of 19,900 connections using all labels before the split raises fold-mean AUROC from 0.756 to 0.987, whereas global selection of 10--100 of 196 regions changes AUROC by only about 0.006 on average. This analysis is a methodological check rather than a study contribution, and the two comparisons differ in cohort, shrinkage, and ranking statistic, so they do not isolate feature granularity alone. Classifier choice also matters. On identical static tangent features and folds, the five evaluated linear classifiers reach AUROCs of 0.735--0.757, whereas the three fixed-configuration tree ensembles reach 0.611--0.624. Logistic regression exceeds each tree ensemble with unadjusted DeLong $p\le2.4\times10^{-12}$. Because only logistic regression is tuned, this comparison applies to the evaluated configurations rather than to the best attainable model from each family, although the other fixed linear classifiers also outperform the tree ensembles. The symptom analyses clarify what the classifier score does not measure. 

Across both diagnostic groups, 3 of 6 score--symptom associations survive correction, whereas none survives within autistic participants. For SRS total, the correlation falls from 0.438 across both groups to 0.050 within autistic participants. The pooled associations therefore largely reflect diagnostic-group separation rather than variation in symptom severity within autism. The windowed-connectivity analysis provides a similar caution about interpretation. The windowed representation initially exceeds the unshrunk static baseline by 0.010 AUROC ($p=0.281$), but applying the same additional covariance shrinkage of 0.15 to the static representation raises its AUROC by 0.013 ($p=0.004$). The windowed representation then differs from the matched static baseline by only $-0.002$ ($p=0.709$). No additional benefit of windowing is therefore detected beyond matched regularisation (difference $-0.002$, 95\% site-clustered interval $-0.015$ to $+0.009$), which shows that static and dynamic representations should be compared under matched shrinkage but does not show that temporal information is never useful. Calibration analysis adds a further distinction between discrimination and probability quality. The native probabilities have slopes close to one but show a small average overprediction of autism. This offset is consistent with penalisation of the intercept as a major contributor, although the unpenalised refit overshoots slightly and does not isolate a single cause, and it is removed by an intercept-only recalibration fitted within each training partition. The correction restores calibration-in-the-large without materially changing AUROC, although it shifts the balance between sensitivity and specificity at the fixed threshold of 0.5. Calibration also remains heterogeneous across sites under LOSO validation, so the correction improves average calibration in this cohort but does not establish calibration at a new site. The attribution analysis shows that interpretation is more stable at the level of functional systems than at the level of individual connections. Default-mode pairs account for the largest overall shares of classifier separation, partly because they contain many connections, while the somatomotor system contributes most per connection. These patterns agree with the adjusted group-difference analysis, particularly for the correlation view. System-level structure coefficients also reproduce across non-overlapping halves of the cohort, which contain the same sites, whereas individual tangent-space connections are less stable. Removing the highest-ranked system pairs lowers held-out AUROC more than random removal, confirming that these pairs contribute to the fitted model, although information remains distributed across many systems. Finally, the secondary analyses in the Supplementary Material further limit the claims supported by these data: per-site abstention answers only 19.5\% of participants and gives no broadly useful guarantee, combining parcellations does not improve discrimination, and the accuracy attainable against the diagnostic labels cannot be identified from two instruments applied to one population. Overall, the results show that performance, calibration, and interpretation depend not only on the connectivity signal but also on cohort composition, regularisation, feature-selection scope, estimator design, and the level at which model explanations are examined.

\subsection{Limitations}
\label{subsec:limitations}

Several considerations define the scope of these findings. The primary cohort spans 6.5 to 58 years, with a median age of 14.8 years, and contains more male than female participants. The results therefore represent the demographic composition of ABIDE-I and should not be extended directly to early identification or more balanced clinical populations. The comparison group is predominantly typically developing rather than clinically referred, so the reported discrimination reflects separation within a research cohort rather than differential diagnosis. LOSO validation provides an important test of transfer across the acquisition sites represented in ABIDE-I, but independent datasets will be needed to determine how well the findings generalise across populations, scanners, and preprocessing pipelines. The functional-connectivity contrasts are also conditional on the selected preprocessing strategy, CC200 parcellation, nuisance model, and the absence of global signal regression. Motion-matched and regional analyses are consequently interpreted as exploratory, while agreement between group-level effects and classifier patterns is examined descriptively rather than as evidence of causation. Classifier comparisons should likewise be read in the context of the evaluated configurations, because logistic regression is tuned whereas the tree ensembles use fixed settings. Additional sensitivity analyses, including alternative atlases, harmonisation, demographic information, and reduced region sets, do not materially improve performance and help define the conditions under which the main results are obtained. The calibration and attribution analyses are assigned descriptive, secondary, or exploratory status according to their analysis timing and purpose. Site-level calibration estimates are specific to the observed sites, and intervals for recalibrated probabilities are conditional on the fitted correction. Structure coefficients may also reflect participant or site characteristics that are not explicitly removed from the classifier, and connections involving four regions with intermittent constant signals warrant cautious interpretation. Reproducibility is assessed using non-overlapping cohort halves, which provides a stricter test than overlapping folds but uses fewer participants in each fit; both halves contain participants from every site, so the test is within the cohort and its site mix and does not show reproducibility across sites or in an independent cohort. Overall, these considerations mainly constrain the generalisability and biological interpretation of the findings rather than the internal consistency of the observed results. The somatomotor and default-mode patterns are therefore presented as cohort-level associations that motivate independent validation rather than as established biomarkers.

\subsection{Assessment against criteria for trustworthy classification}
\label{subsec:trustworthiness}

Trustworthiness is not a single property, so the evidence of this study is set against separate criteria in Table~\ref{tab:trustworthiness}. The ensemble meets the criteria for average probability calibration after training-only recalibration and for system-level reproducibility of its explanations within the cohort. It does not meet, or the study cannot establish, the criteria for calibration at a new site, a usable operating point, independence from cohort composition, and external validity. The study therefore provides a framework for assessing trustworthiness together with evidence on one ensemble, but it does not establish that this classifier is trustworthy.

\begin{table}[t]
\centering
\caption{Evidence of this study against criteria for trustworthy connectome-based autism classification.}
\label{tab:trustworthiness}
\small
\begin{tabular}{@{}>{\raggedright\arraybackslash}p{0.22\linewidth}>{\raggedright\arraybackslash}p{0.52\linewidth}>{\raggedright\arraybackslash}p{0.20\linewidth}@{}}
\toprule
Criterion & Evidence & Status \\
\midrule
Discrimination at an unseen site & LOSO AUROC 0.774; site AUROC 0.589--0.912; the intervals of 5 of 20 sites include 0.5 & Partly met \\
Average probability calibration & Calibration-in-the-large $-0.162$ to $-0.010$ after training-only recalibration; AUROC unchanged & Met within the cohort \\
Calibration at a new site & After correction, LOSO site-level heterogeneity $I^2=63\%$; the intervals of two sites exclude zero & Not established \\
Usable operating point & Recalibration changes sensitivity from 0.705 to 0.655 (pooled) and from 0.742 to 0.625 (LOSO) at 0.5; abstention gives no useful guarantee & Not established \\
Independence from cohort composition & AUROC 0.873 to 0.719 across unmatched IQ bands; stronger score--IQ association in autistic participants for full-scale and verbal IQ & Not met \\
Reproducible explanations & System-pair structure coefficients 0.60--0.70 between halves against permutation 95th percentiles of at most 0.13; single connections 0.18--0.53 & Met at system level, within the cohort \\
External validity & One research cohort (ABIDE-I); no independent dataset & Not established \\
\bottomrule
\end{tabular}
\end{table}

\subsection{Additional analyses and remaining requirements}
\label{subsec:required-validation}

Additional analyses clarify the uncertainty estimates, nesting of the design, and cognitive-ability findings. Site-clustered bootstrap intervals are wider than the DeLong and Wilson intervals, particularly at their lower bounds (Section~\ref{subsec:results-performance}), indicating that the parametric intervals understate some of the uncertainty associated with site clustering. However, all AUROC intervals remain above 0.5. Calibration is also evaluated using several measures rather than ECE alone (Section~\ref{subsec:results-calibration}), because the binned error does not fully capture the average probability offset. The cognitive-ability result is tested directly through diagnosis-by-IQ interactions. These interactions are negative under both validation protocols and survive correction for full-scale and verbal IQ, but not for performance IQ (Section~\ref{subsec:results-cognitive-ability}). The interpretation is therefore limited to full-scale and verbal IQ, and all analysed cohort counts are reconciled with explicit inclusion rules from the released records. The main results are also insensitive to stricter nesting. Refitting the tangent reference and feature scaler within every inner split changes the selected penalty in only 2 of 30 member fits under pooled validation and 2 of 60 under LOSO validation, while participant-pooled AUROC changes by less than 0.001 (Section~\ref{subsec:results-performance}). Similarly, refitting the nuisance model using only outer-training participants changes the region-selection AUROC by 0.0035 on average (Section~\ref{subsec:results-selection-scope}). The group-difference tests use Freedman--Lane residual permutation, and the within-site permutation of residualised connections is reported as a sensitivity analysis that identifies a subset of the same pairs (Section~\ref{subsec:results-system-differences}). Independent external validation nevertheless remains necessary. The study is therefore interpreted as a retrospective analysis of an aggregated research cohort, without claims about clinical screening, deployment, or performance at a new acquisition site. Its screening-related implications are methodological and descriptive, concerning how classifier scores are calibrated, interpreted, and reported rather than supporting their clinical use.

\subsection{Future work}
\label{subsec:future-work}

Future work will extend the present findings toward broader validation, stronger generalisability, and more consistent reporting of connectome-based autism classification. External validation will be conducted on an independent cohort such as ABIDE-II, either by reproducing the present preprocessing pipeline or by explicitly evaluating transfer across preprocessing strategies. The influence of cognitive composition will also be examined in cohorts with more balanced IQ distributions, greater representation of female participants and younger children, and clinically referred comparison groups. The somatomotor and default-mode findings will be reassessed under alternative denoising strategies, including global signal regression and more detailed motion modelling, while the contribution of regions affected by constant signals will be examined separately. Classifier comparisons will be expanded to tuned nonlinear models under the same nested validation framework, allowing model complexity to be evaluated more fairly. Calibration and system-level attribution will also be tested across different solvers, preprocessing pipelines, and independent datasets to determine how consistently these properties are retained. More broadly, future studies will benefit from reporting discrimination together with calibration, operating thresholds, cohort characteristics, and the reproducibility of model explanations. This direction will support more transparent comparison across studies and will help establish which aspects of connectome-based autism classification remain stable across populations, sites, and analytical settings.